% Ad Atticum 6.4 (ad-atticum-06-04) --- English translation
% Translated by Claude (Anthropic) from the Latin (Perseus).
% Style guide: STYLE.md.

\ciceroLetterOpener{CICERO ATTICO salutem}

\ciceroSection{1} We came to Tarsus on the Nones of June. Many things weighed on me there: a great war in Syria, great brigandage in Cilicia, and for me a difficult task of administration --- because I had only a few days left of my year of office, and what was hardest of all: by the senate's decree someone had to be left in charge. Nothing could have been less approved than the quaestor Mescinius. As for Coelius, we were hearing nothing. It seemed most correct to leave my brother with the command; but in that there were many things to grieve --- our parting, the danger of the war, the unruliness of the soldiers, six hundred other things besides. What a hateful business altogether! But let fortune look to all this, since deliberation is not much allowed me.

\ciceroSection{2} You, then, since by now you have come safe to Rome, as I hope, will see to it, as you always do, in everything that you know touches my interest --- above all about my Tullia, on whose match I wrote what I thought to Terentia while you were in Greece; and then about my own honour. For since you were away, I am afraid the matter of my despatch was not handled in the Senate with sufficient diligence. There is one further thing I shall write to you about \textit{[Greek: more in the mystery-key]}; you, more keenly, will sniff it out. \textit{[Greek: My wife's freedman --- you know the one --- it seemed to me the other day, from what he was letting fall in his calculations, that he has been juggling the accounts in the matter of the purchase of the holdings of\ldots that you may understand. One man, presumably.]} I cannot write as much as I fear; but you --- see to it that your letters come flying to meet me. I have written this in haste, on the road and on the march. Give my best greetings to Pilia and the lovely little Caecilia.
